\documentclass[11pt,a4paper,twoside]{article}

\usepackage[T1]{fontenc}
\usepackage[utf8]{inputenc}
\usepackage[spanish,es-tabla,es-nodecimaldot,es-noshorthands]{babel}
\spanishplainpercent
\usepackage{lmodern}
\usepackage{microtype}
\usepackage[a4paper,inner=23mm,outer=23mm,top=23mm,bottom=23mm,headheight=14pt,headsep=8mm]{geometry}
\usepackage{setspace}
\setstretch{1.10}
\raggedbottom
\usepackage{amsmath,amssymb,mathtools}
\usepackage{graphicx}
\graphicspath{{figuras/}}
\usepackage[table]{xcolor}
\definecolor{azul}{HTML}{0072B2}
\definecolor{optimista}{HTML}{009E73}
\definecolor{pesimista}{HTML}{D55E00}
\definecolor{tinta}{HTML}{252A30}
\definecolor{gris}{HTML}{60666D}
\definecolor{fondo}{HTML}{F1F5F8}
\color{tinta}
\usepackage{booktabs,tabularx,longtable,array,makecell,multirow}
\renewcommand{\arraystretch}{1.16}
\newcolumntype{Y}{>{\raggedright\arraybackslash}X}
\usepackage{siunitx}
\sisetup{output-decimal-marker={,},group-separator={.},group-minimum-digits=4}
\usepackage{eurosym}
\usepackage{enumitem}
\setlist{nosep,leftmargin=*,topsep=4pt}
\usepackage{caption,subcaption}
\captionsetup{font=small,labelfont=bf,labelsep=period,justification=justified,singlelinecheck=false}
\captionsetup[table]{position=top,skip=6pt}
\usepackage{float,placeins,pdflscape}
\usepackage{titlesec}
\titleformat{\section}{\Large\bfseries}{\thesection.}{.65em}{}
\titleformat{\subsection}{\large\bfseries}{\thesubsection.}{.65em}{}
\titleformat{\subsubsection}{\normalsize\bfseries}{\thesubsubsection.}{.65em}{}
\titlespacing*{\section}{0pt}{20pt plus 3pt minus 2pt}{9pt}
\titlespacing*{\subsection}{0pt}{14pt plus 2pt}{6pt}
\usepackage{fancyhdr}
\pagestyle{fancy}
\fancyhf{}
\fancyhead[LE,RO]{\small\thepage}
\fancyhead[LO]{\small\color{gris}Kaido Motors España · Guía del proceso}
\fancyhead[RE]{\small\color{gris}\nouppercase{\leftmark}}
\renewcommand{\headrulewidth}{0pt}
\renewcommand{\sectionmark}[1]{\markboth{\thesection.\ #1}{}}
\usepackage[skins]{tcolorbox}
\newtcolorbox{idea}[1]{enhanced,colback=fondo,colframe=gris!35,
  boxrule=.4pt,arc=0pt,left=9pt,right=9pt,top=7pt,bottom=7pt,
  title={#1},colbacktitle=fondo,coltitle=tinta,fonttitle=\bfseries,
  titlerule=0pt,before skip=10pt,after skip=10pt}
\usepackage{tikz}
\usetikzlibrary{arrows.meta,positioning,calc}
\usepackage{pgfplots,pgfplotstable}
\pgfplotsset{compat=1.18,every axis/.append style={
  font=\small,axis line style={gris!60},tick style={gris!60},
  grid=major,grid style={gris!14},axis lines=left,
  legend style={draw=none,fill=white,font=\small},
  label style={font=\small},title style={font=\normalsize\bfseries}}}
\usepackage{xurl}
\usepackage[unicode,colorlinks=true,linkcolor=azul,citecolor=azul,urlcolor=azul,
  pdftitle={Kaido Motors España: guía razonada del proceso y los resultados},
  pdfauthor={Proyecto Kaido Motors},pdfsubject={Simulación de Sistemas Logísticos}]{hyperref}
\usepackage{bookmark}
\urlstyle{tt}
\newcommand{\archivo}[1]{\nolinkurl{#1}}
\newcommand{\fuente}[1]{\par\smallskip{\footnotesize\color{gris}\textit{Fuente:} #1.\par}}
\newcommand{\unidad}[1]{\,\text{#1}}
\newcommand{\MEE}{\text{M\euro}}
\newcommand{\paso}[1]{\paragraph{#1}}
\newcommand{\tabla}[4]{%
  \begin{table}[htbp]\centering\caption{#2}\label{#3}
  {\small\input{tablas/#1}}\fuente{#4}\end{table}}
\newcommand{\grafica}[4]{%
  \begin{figure}[htbp]\centering
  \includegraphics[width=\linewidth,height=.65\textheight,keepaspectratio]{#1}
  \caption{#2}\label{#3}\fuente{#4}\end{figure}}
\newcommand{\grafosanchos}[5][]{%
  \clearpage\begin{landscape}\thispagestyle{plain}
  \begin{figure}[p]\centering\includegraphics[width=.96\linewidth]{#2}
  \caption{#3}\label{#4}\fuente{#5}
  \if\relax\detokenize{#1}\relax\else
    \medskip\begin{minipage}{.94\linewidth}\small #1\end{minipage}
  \fi\end{figure}
  \end{landscape}\clearpage}
\setcounter{tocdepth}{2}
\setlength{\parindent}{1.25em}
\setlength{\parskip}{3pt plus 1pt}
\setlength{\emergencystretch}{2em}
\widowpenalty=10000
\clubpenalty=10000
\renewcommand{\topfraction}{.88}
\renewcommand{\bottomfraction}{.8}
\renewcommand{\textfraction}{.1}
\renewcommand{\floatpagefraction}{.7}

\begin{document}
\thispagestyle{empty}
\begin{center}
  \vspace*{4mm}
  \includegraphics[width=4.2cm]{logo-uax.png}\par
  \vspace{9mm}
  {\scshape Simulación de Sistemas Logísticos}\par
  \vspace{6mm}
  {\LARGE\bfseries Kaido Motors España}\par
  \vspace{3mm}
  {\Large Demanda, stock y beneficio}\par
  \vspace{7mm}
  {\large\color{azul}Guía razonada del proceso y los resultados}\par
  \vspace{7mm}
  {Universidad Alfonso X el Sabio}\par
  {\small Ejecución de referencia: 10 de octubre de 2026}\par
\end{center}
\vspace{7mm}

\noindent\textbf{Resumen.}
Este informe explica cómo se transforma el histórico de Kaido Motors en una
propuesta de pedidos para A013. Presenta la preparación de datos, la separación
entre matriculaciones y compras de clientes, la previsión de mercado, la
simulación mensual con stock de distintas edades y la elección de una política
de abastecimiento. Cada fase conecta la pregunta de negocio con su método,
sus resultados y sus límites. Las tablas proceden de la ejecución revisada de
los cuatro notebooks del proyecto; las cifras de escenarios son condicionales
a sus supuestos y no tienen probabilidades asignadas.

\medskip
\noindent\textbf{Palabras clave:} demanda de clientes · validación temporal ·
stock de seguridad · FIFO · Monte Carlo · contribución económica.

\begin{idea}{Cómo utilizar este documento}
Para comprender el conjunto, leer las secciones~\ref{sec:mapa} y
\ref{sec:lectura-final}. Para reproducir el razonamiento, seguir las cuatro
fases en orden. Las ecuaciones incluyen la interpretación de sus símbolos;
los ejemplos pequeños permiten comprobar una cuenta antes de pasar a las
tablas reales. Los anexos reúnen planes mensuales, resultados por segmento y
preguntas de repaso.
\end{idea}

\vfill
\noindent\footnotesize
\textbf{Referencia docente.} Caso original elaborado por Jacinto Velasco y
Javier Goikoetxea. Profesor: Javier Goikoetxea. Este documento es una explicación
de la solución desarrollada en el proyecto; no es el enunciado ni se atribuye su
autoría a los autores del caso. Los datos y nombres están camuflados. Material de
uso exclusivamente docente, sujeto a las condiciones del enunciado.
\normalsize
\clearpage
\tableofcontents
\clearpage

\section{El problema completo, antes de entrar en los cálculos}\label{sec:mapa}

\subsection{La decisión que queremos tomar}
Kaido debe decidir cuánto encargar a fábrica sin conocer las ventas futuras.
Si encarga poco, puede perder compradores; si encarga demasiado, mantiene coches
que cuestan dinero y envejecen. Además, el pedido no llega de inmediato. Una
decisión tomada después de observar una subida de ventas tarda tres meses en
convertirse en una llegada modificable a la red.

El objetivo del trabajo no es producir una cifra de ventas aislada. Es construir
una \textbf{regla de decisión que siga funcionando cuando la previsión se
equivoca}. Por eso se necesitan dos herramientas diferentes:
\begin{enumerate}
  \item Una \textbf{previsión}: estima cuántos clientes querrán comprar cada mes.
  \item Una \textbf{simulación}: representa qué puede venderse realmente con el
  stock disponible, los compromisos de suministro y la regla de revisión.
\end{enumerate}
La recomendación económica se obtiene al comparar políticas sobre las mismas
trayectorias de demanda. El modelo trabaja por mes y por segmento; no representa
individualmente los 64 concesionarios.

\begin{figure}[htbp]
\centering
\begin{tikzpicture}[node distance=5mm,
  caja/.style={draw=gris!40,fill=fondo,rounded corners=1pt,text width=3.2cm,
    minimum height=1.65cm,align=center,font=\small},
  flecha/.style={-{Stealth[length=2mm]},draw=azul,thick}]
  \node[caja] (d) {\textbf{Datos}\\Qué ocurrió\\Clientes y episodios};
  \node[caja,right=of d] (p) {\textbf{Previsión}\\Qué podría ocurrir\\Mercado y posición};
  \node[caja,right=of p] (s) {\textbf{Simulación}\\Qué podremos atender\\Stock y plazos};
  \node[caja,right=of s] (e) {\textbf{Decisión}\\Qué política elegir\\Servicio y margen};
  \draw[flecha] (d)--(p);\draw[flecha] (p)--(s);\draw[flecha] (s)--(e);
\end{tikzpicture}
\caption{Las cuatro piezas de la solución. Un error en una fase se transmite a
las siguientes: una demanda mal interpretada puede terminar en un pedido mal dimensionado.}
\label{fig:proceso}
\fuente{estructura de los cuatro notebooks revisados}
\end{figure}

\subsection{Cómo se organiza el proyecto}
\begin{table}[htbp]\centering
\caption{Responsabilidad y resultado de cada notebook.}\label{tab:mapa-notebooks}
\begin{tabularx}{\linewidth}{@{}p{3.4cm}YY@{}}\toprule
Notebook & Pregunta principal & Resultado que entrega \\\midrule
\archivo{datos.ipynb} & ¿Qué parte del histórico representa compras de clientes?
& Series observadas, diagnóstico y ajustes de episodios.\\
\archivo{demanda.ipynb} & ¿Qué demanda mensual podemos prever y qué error tiene?
& Previsión, evaluación temporal y parámetros de incertidumbre.\\
\archivo{simulacion.ipynb} & ¿Cómo se transforman pedidos y llegadas en ventas y costes?
& Simulador comprobado y referencia determinista.\\
\archivo{escenarios.ipynb} & ¿Qué política compensa y qué riesgos mantiene?
& Simulaciones, política elegida, sensibilidades y entrega.\\\bottomrule
\end{tabularx}
\end{table}

Los notebooks comparten las celdas de definiciones etiquetadas \archivo{api}.
Las fases posteriores cargan esas funciones en un espacio independiente; no
ejecutan de nuevo los gráficos de otra fase ni usan las tablas antiguas como si
fueran el modelo. \archivo{analisis.ipynb} y los módulos históricos de previsión
y simulación se conservan como referencia inicial. Este informe describe la
\textbf{versión revisada}.

\subsection{Tres niveles de afirmación}
Conviene identificar en cada resultado de dónde procede la cifra:
\begin{description}[style=nextline]
  \item[Observado.] Se lee en los datos o se calcula directamente a partir de
  ellos: por ejemplo, matriculaciones menos automatriculaciones.
  \item[Estimado o supuesto.] Requiere un método o una hipótesis: por ejemplo,
  cuántas compras se adelantaron o cuántos coches vienen ya en camino.
  \item[Simulado.] Es el resultado de aplicar reglas y supuestos a una posible
  trayectoria futura: ventas, stock y contribución de una réplica.
\end{description}
La distinción importa porque repetir mil veces una simulación no convierte sus
supuestos en hechos. Tampoco una correlación histórica demuestra que una variable
sea la causa de un cambio en las ventas.

\subsection{Vocabulario mínimo para leer el resto}
\begin{table}[htbp]\centering
\caption{Conceptos que no deben confundirse.}\label{tab:glosario-inicial}
\begin{tabularx}{\linewidth}{@{}p{3.2cm}Y@{}}\toprule
Concepto & Interpretación en este trabajo \\\midrule
Matriculaciones & Incluyen ventas a clientes y automatriculaciones.\\
Venta a cliente & Compra observada después de descontar automatriculaciones.\\
Demanda potencial & Lo que los clientes querrían comprar; puede superar las ventas si falta stock.\\
Wholesale / retail & Facturación a la red / venta al cliente, respectivamente.\\
PFF & Precio franco fábrica; no es el PVP con impuestos.\\
VM & Gasto comercial: descuentos, campañas e incentivos por vehículo.\\
CVP & Contribución del vehículo después del VM: beneficio bruto menos VM.\\
Transfer & Vehículo procedente del stock de otro país europeo para cubrir un faltante.\\
FIFO & Se vende primero la cohorte más antigua disponible.\\
Escenario & Conjunto coherente de supuestos sobre A013.\\
Réplica & Una trayectoria aleatoria completa dentro de un escenario.\\\bottomrule
\end{tabularx}
\fuente{enunciado, secciones 2 y 9, y convenciones de la simulación revisada}
\end{table}

\begin{idea}{La idea que conviene conservar}
Una previsión de \num{1000} coches no significa que se venderán \num{1000}.
Primero tiene que existir esa demanda y después tiene que haber coches para
atenderla. Por último, vender esos coches debe evaluarse con sus costes: más
unidades no siempre implican más contribución.
\end{idea}
\FloatBarrier

\section{Primera fase: entender y preparar los datos}\label{sec:datos}

\begin{idea}{Entrada, trabajo y salida}
\textbf{Entran} siete CSV del profesor. \textbf{Se comprueban} claves, periodos,
ausencias y balances; se separan clientes y automatriculaciones y se estudian los
episodios excepcionales. \textbf{Salen} una serie mensual por segmento, un
diagnóstico de A012 y alternativas explícitas para lo que no se observa.
\end{idea}

\subsection{Qué contiene cada fuente}
Los ficheros no empiezan todos en el mismo año. Mercado y economía comienzan en
A001; las ventas internas, en A003; el stock, en A009. Una fila de ventas tiene
tres claves: mes, segmento y canal. Esto explica las \num{2520} filas:
$10\text{ años}\times12\text{ meses}\times7\text{ segmentos}\times3\text{ canales}$.
La tabla de stock no tiene canal, de modo que no se puede cruzar con ventas como
si cada canal dispusiera de su propio inventario.

\tabla{inventario.tex}{Datos originales disponibles y granularidad de cada fila.}
{tab:inventario}{\archivo{data/README.md} y anexo 5 del enunciado}

Los años A001--A013 son códigos del caso, no años naturales identificables. El
contador $t$ permite ordenar los meses: diciembre de A012 es $t=144$ y enero de
A013 es $t=145$. Las fechas sintéticas que utiliza una biblioteca no proporcionan
el calendario laboral: días laborables y Semana Santa se leen en los CSV.

\subsection{Por qué limpiar no significa cambiar los datos}
La lectura preserva el código de modelo \archivo{NA}; no lo convierte en un
valor ausente. Un campo vacío sí representa información no disponible. Un cero,
en cambio, es una cantidad observada: por ejemplo, ventas nulas antes de un
lanzamiento. Las columnas porcentuales tampoco comparten siempre escala:
\archivo{pct_diesel} se expresa como proporción, mientras el paro se expresa en
puntos de porcentaje.

Las comprobaciones verifican que las claves no estén duplicadas, que los periodos
sean coherentes y que las sumas y desgloses tengan sentido. Pasan los 25 controles
de lectura y consistencia. Esto \textbf{no implica que todas las identidades
históricas sean exactas}: las discrepancias se muestran sin alterar los originales.

\paso{El balance histórico.}
Si $S_t$ es el stock al cerrar el mes, $L_t$ la facturación a la red y $M_t$ las
matriculaciones, el residuo es
\begin{equation}
 r_t=S_t-\bigl(S_{t-1}+L_t-M_t\bigr).
\end{equation}
En el histórico existen diferencias pequeñas y algunos casos mayores: el máximo
absoluto es de 51 coches en SSUV y de 8 en MD. En MSUV el máximo es 5; en CSUV,
cero. Se mantienen las discrepancias, con localización en las tablas de diagnóstico.
No se fuerza el residuo a cero para mejorar el aspecto de los datos.

\fuente{\archivo{resultados/datos/tablas/balance_stock_segmento.csv}}

Es importante usar aquí matriculaciones y no sustituirlas por ventas netas. Las
automatriculaciones forman parte de las salidas registradas en el balance
histórico. Si se retiran sólo de las salidas y no se cambia el sistema de
contabilidad del stock, se introduce un descuadre artificial.

\subsection{De matriculaciones a compras observadas de clientes}
Para cada mes y segmento se define
\begin{equation}\label{eq:neto}
 O_{s,t}=M_{s,t}-A_{s,t},
\end{equation}
donde $O$ son las compras observadas de clientes, $M$ las matriculaciones y $A$
las automatriculaciones; todas las cantidades se expresan en coches. Las ventas
especiales no se restan automáticamente: una venta a una empresa puede ser una
compra real de un cliente.

En A012 se registran \num{23175} matriculaciones y \num{2469}
automatriculaciones. Por tanto, se observan \num{20706} compras de clientes.
En A011 se observan \num{19543}. La variación es
\begin{equation}
 g_{\mathrm{observado}}=
 \frac{\num{20706}}{\num{19543}}-1\simeq 6{,}0\,\%.
\end{equation}
La serie resultante sigue sin ser toda la demanda potencial. Puede incluir
compras adelantadas y puede estar limitada por falta de coches.

\subsection{El lanzamiento de CSUV cambia la comparación}
CSUV empieza a venderse en agosto de A012 y aporta \num{1192} clientes ese año.
Son clientes reales y pertenecen al total de Kaido. Al excluirlos queda una gama
comparable con \num{19514} clientes, prácticamente estable frente a A011:
$-0{,}1\,\%$.

No hay contradicción entre que crezca la marca y no crezca su gama anterior. Son
dos preguntas distintas: cuánto vende el conjunto y cómo evoluciona el perímetro
de productos que ya existía.

\tabla{crecimiento.tex}{Crecimiento observado y ajuste parcial de episodios.
Los valores ajustados son estimaciones, no recuentos de clientes identificados.}
{tab:crecimiento}{\archivo{resultados/datos/tablas/crecimiento_demanda_a012.csv}}

\grafica{crecimiento_demanda_observado_ajustado_a012.png}
{A012 frente a A011: la conclusión depende de mantener el perímetro de comparación.
El total incluye CSUV; el ajuste central trata episodios, pero no identifica toda
la demanda potencial.}
{fig:crecimiento}{\archivo{datos.ipynb}, gráfico de crecimiento observado y ajustado}

\subsection{Cómo se identifican las compras adelantadas}
Una compra adelantada cambia de fecha; no crea necesariamente un cliente nuevo.
Los indicadores del caso señalan un episodio en agosto de A011, seguido de caída
en septiembre--diciembre, y otro en noviembre--diciembre de A012.

Para estimar el exceso se utilizan A009 y A010 como referencias estacionales.
Cada referencia se escala \emph{por separado} al nivel observado antes del
episodio del año que se estudia. Para una referencia $r$, un segmento $s$ y un
año $a$, la cuenta conceptual es
\begin{align}
 k_{s,a,r}&=\frac{\text{clientes antes del episodio en }a}
                  {\text{clientes en los mismos meses de }r},\\
 E_{s,a,r}&=\max\!\left(0,\ \text{clientes en meses del episodio de }a
                 -k_{s,a,r}\,\text{clientes en esos meses de }r\right).
\end{align}
El mínimo y máximo proceden de esas dos referencias y el central es su punto
medio. El recorte a cero evita atribuir un adelanto negativo; también recuerda que
la estimación es una regla de identificación, no una observación directa.

\tabla{adelantos.tex}{Exceso estimado en los episodios y compensación posterior.
La amplitud es sensibilidad a referencias, no un intervalo estadístico.}
{tab:adelantos}{\archivo{resultados/datos/tablas/adelantos_a011_a012.csv}}

\paso{A011: estudiar el pico y la caída conjuntamente.}
Se estiman \num{450.4} coches de exceso y \num{257.3} de caída posterior
compensatoria, limitada al exceso del segmento. La estimación anual ajustada es
\begin{equation}
 \num{19543}-\num{450.4}+\num{257.3}\simeq\num{19349.9}.
\end{equation}
El saldo no compensado no demuestra que esas compras se tomaran de A012: el
episodio fue en agosto y la referencia estacional es imperfecta. Por eso no se
convierte automáticamente en una transferencia cierta al año siguiente.

\paso{A012: registrar una sola devolución hacia A013.}
El exceso central de cierre se estima en \num{878.4} coches. Al restarlo una
sola vez se obtienen \num{19827.6} compras ajustadas en A012. Comparadas con
las \num{19349.9} ajustadas de A011, implican un aumento parcial del 2,5\,\%.

El registro mensual reparte el ajuste anual entre los meses del episodio. No
resta el exceso anual completo en cada mes. La misma magnitud alimentará después
la caída supuesta en A013: retirar la anomalía del histórico y representar su
devolución futura cumplen funciones diferentes, pero no deben generar dos
descuentos sobre la misma previsión futura.

\paso{Por qué CSUV necesita otro tratamiento.}
No existe una historia de ventas propias anterior a su lanzamiento. Se utiliza
como proxy el exceso de su mercado en noviembre--diciembre respecto a referencias
escaladas, aplicado a las compras de clientes del segmento. En esta ejecución no
se detecta un exceso positivo con ese método. Es un resultado del proxy, no una
demostración de que ningún comprador adelantara su compra.

\subsection{Qué sabemos y qué no sabemos sobre la falta de stock}
Una señal de stock libre bajo identifica un mes sospechoso, pero no cuenta a
los clientes que se marcharon. Si se disponía de diez coches y se vendieron
diez, la demanda pudo ser diez o veinte. El volumen perdido no se observa en
esas ventas.

El indicador de preparación compara el stock sin cliente con las
matriculaciones del mes y utiliza medio mes de cobertura como umbral. No es el
nivel de servicio futuro ni el stock de seguridad: sirve para señalar un posible
problema de observación del histórico.

Como sensibilidad se utiliza la relación entre mercado y compras de clientes
en hasta seis meses anteriores sin señal de restricción, adelanto ni caída.
Si esa referencia sugiere una demanda mayor, se construye
\begin{equation}
 O^{\mathrm{alternativa}}_{s,t}
 =\max\!\left(O_{s,t},\frac{I_{s,t}}
 {\operatorname{mediana}_{u<t}\bigl(I_{s,u}/O_{s,u}\bigr)}\right),
\end{equation}
con $I$ igual a las matriculaciones del mercado y usando sólo antecedentes
válidos. Si no hay referencia suficiente, se conserva lo observado.

La alternativa arroja un crecimiento observado ampliado del 6,8\,\% para el
total, frente al 6,0\,\% de las compras registradas. No es una estimación
independiente de la demanda perdida: depende de la estabilidad de esa relación
comercial. La previsión principal se valida contra clientes \emph{observados},
no contra una serie inventada por el propio modelo.

\grafica{stock_red_antiguedad_a011_a012.png}
{Stock observado de la red por edad. La cantidad total y la antigüedad son dos
dimensiones distintas: reducir coches no garantiza que desaparezca primero el
inventario problemático.}
{fig:stock-historico}{\archivo{datos.ipynb}; stock mensual A011--A012}

\begin{idea}{Qué concluye realmente la primera fase}
Se observa crecimiento de clientes en el total y estabilidad de la gama anterior.
El ajuste central de episodios sigue siendo positivo para el total, pero depende
de las referencias y de la compensación de A011. La señal de stock y la historia
breve de CSUV impiden presentar una cifra única como demanda potencial exacta.
La política futura de cero automatriculaciones es un supuesto de negocio, no una
corrección retroactiva de lo ocurrido.
\end{idea}
\FloatBarrier
\grafosanchos{demanda_observada_segmento_a003_a012.png}
{Compras observadas por segmento, A003--A012. Los lanzamientos y los episodios
marcados impiden tratar todas las series como una historia homogénea. Las escalas
son distintas para ver los patrones; no deben compararse alturas entre paneles.}
{fig:historico-clientes}{\archivo{datos.ipynb}; serie mensual neta de automatriculaciones}

\section{Segunda fase: prever la demanda y medir el error}\label{sec:prevision}

\begin{idea}{Entrada, trabajo y salida}
\textbf{Entran} clientes observados, mercado completo, economía y calendario.
\textbf{Se comparan} cuatro modelos con cortes temporales; se revisan sesgos y
se trata CSUV por separado. \textbf{Salen} doce previsiones por segmento,
diagnósticos de error y parámetros que alimentan el stock de seguridad y la
simulación aleatoria.
\end{idea}

\subsection{Por qué separar mercado y posición de Kaido}
El punto de partida es
\begin{equation}\label{eq:mercado-cuota}
 \widehat D_{s,t}=\widehat I_{s,t}\,q_{s,t},
\end{equation}
donde $\widehat I$ es el mercado previsto, en matriculaciones, y $q$ una relación
comercial de clientes de Kaido sobre ese mercado. Si el mercado sube pero Kaido
pierde posición, los dos efectos pueden compensarse. Si ambos mejoran, la demanda
propia puede crecer con fuerza.

$q$ no es exactamente una cuota homogénea: en el numerador se han retirado las
automatriculaciones de Kaido, pero no conocemos todas las de sus competidores.
La relación sirve para predecir; su definición debe conservarse al estimar efectos
de producto y al interpretar los resultados.

La previsión utiliza el mercado desde A001, aunque la venta interna de Kaido
comience en A003. Recortar las dos fuentes al mismo inicio perdería información
de mercado sin una necesidad metodológica.

\subsection{Los cuatro candidatos y lo que aprende cada uno}
Se empieza con una referencia que cualquiera puede calcular: repetir el mismo
mes del año anterior. En el ancla histórica de años cerrados se normalizan los
episodios conocidos, retirando su exceso y restituyendo la caída estimada. Así no
se reproduce automáticamente el pico de noviembre--diciembre de A012 en A013.
Durante una revisión dentro del año no se utilizan ajustes calculados con el
cierre que todavía no se conoce.

Se comparan después tres modelos de mercado:
\begin{enumerate}
  \item Confianza del consumidor, calendario y sucesos, sin constante.
  \item Las mismas variables, con constante o intercepto.
  \item Confianza, paro y tipo de crédito, además de calendario y sucesos,
  con constante.
\end{enumerate}
El intercepto permite un crecimiento de fondo incluso cuando las variaciones
de los indicadores son cero. No es siempre mejor: puede arrastrar una expansión
pasada a un año de caída. Por eso se compara mediante validación, en lugar de
eliminarlo después de mirar el año que queremos evaluar.

\subsection{La ecuación del mercado, explicada término a término}
El modelo elegido trabaja con variación logarítmica interanual. Su forma es
\begin{align}\label{eq:modelo-mercado}
 y_{s,t}&=\log\!\left(\frac{I_{s,t}}{I_{s,t-12}}\right),\\
 y_{s,t}&=\alpha+\beta_W\,\Delta_{12}\log W_t
  +\beta_S\,\Delta_{12}H_t
  +\beta_C\,\Delta_{12}C_{t-3}
  +\sum_j\gamma_j\,\Delta_{12}E_{j,t}+\varepsilon_{s,t}.
\end{align}
Aquí $W_t$ son días laborables, $H_t$ indica Semana Santa, $C_t$ es confianza y
$E_{j,t}$ identifica incentivo público, adelanto, caída posterior o subida de
IVA. El operador $\Delta_{12}x_t=x_t-x_{t-12}$ compara con el mismo mes del año
anterior. En la fase logística se utilizará $L_t$ para las llegadas de coches;
el calendario y el flujo físico tienen símbolos diferentes.

\begin{itemize}
  \item Comparar enero con enero reduce el efecto de la estacionalidad anual.
  \item El logaritmo permite interpretar cambios de forma proporcional; al
  deshacerlo se multiplica el mercado del año anterior por $\exp(\widehat y)$.
  \item Los sucesos evitan atribuir todos los picos a la economía normal.
  \item El retraso de tres meses intenta representar información disponible
  al decidir, pero no hace conocida toda la economía futura.
\end{itemize}

Se estima una relación común con observaciones de los segmentos. La regresión
ponderada minimiza
\begin{equation}
 \sum_{s,t} w_{s,t}\bigl(y_{s,t}-\widehat y_{s,t}\bigr)^2,
 \qquad w_{s,t}=\sqrt{I_{s,t-12}}.
\end{equation}
Los mercados grandes pesan más, aunque menos que proporcionalmente a su tamaño.
Se excluyen del ajuste variaciones cuya base sea inferior a 200 coches, porque
un denominador muy pequeño produce tasas extremadamente inestables. La regla
no equivale a borrar esos segmentos de la entrega: reciben previsión y diagnóstico.

\subsection{Qué significa respetar la información disponible}
Supongamos que estamos al cerrar A011. El fichero contiene toda la economía
A012, pero no está permitido usarla en una previsión operativa realizada en
esa fecha. Para un indicador retrasado, el origen máximo que puede consultarse es
\begin{equation}
 \text{mes económico utilizado}=\min(t-3,\,t_0),
\end{equation}
siendo $t_0$ el cierre de información. Si el dato necesario todavía no se conoce,
se mantiene el último disponible en el central o se introduce la trayectoria
supuestamente definida para un escenario.

El calendario sí se conoce. Los sucesos futuros de las evaluaciones operativas
se fijan en cero; otra evaluación, claramente separada, utiliza economía y
sucesos reales y se denomina \emph{condicionada}. No puede presentarse como el
rendimiento alcanzable sin conocer el futuro.

\subsection{Cómo se selecciona el modelo sin utilizar A012}
Se hacen tres ejercicios de doce meses:
\begin{center}
\begin{tabular}{@{}lll@{}}\toprule
Información hasta & Periodo previsto & Función \\\midrule
A008, $t_0=96$ & A009 & Selección \\
A009, $t_0=108$ & A010 & Selección \\
A010, $t_0=120$ & A011 & Selección y contraste de sesgo \\
A011, $t_0=132$ & A012 & Evaluación retrospectiva separada \\\bottomrule
\end{tabular}
\end{center}

Para elegir se utiliza WAPE, calculado sobre todos los meses y los seis segmentos
validables:
\begin{equation}\label{eq:wape}
 \mathrm{WAPE}=\frac{\sum_{s,t}|\widehat D_{s,t}-O_{s,t}|}
                         {\sum_{s,t}O_{s,t}}.
\end{equation}
CSUV no participa porque no posee historia comparable. Se elige el candidato con
menor WAPE; la referencia más sencilla tiene prioridad en un empate.

\tabla{seleccion.tex}{Comparación de candidatos con A009--A011, antes de A012.}
{tab:seleccion}{\archivo{resultados/demanda/tablas/seleccion_a009_a011.csv}}

Gana el modelo con confianza e intercepto, con WAPE del 23,8\,\%. No significa
que todos los segmentos tengan ese error ni que sea un error futuro garantizado.
Se trata de una comparación entre cuatro alternativas bajo el protocolo definido.

\subsection{Por qué el sesgo se estudia aparte del error}
El error se define como previsión menos observado. Un sesgo positivo significa
que se prevé demasiado; uno negativo, que se prevé poco. Dos modelos pueden
tener errores absolutos parecidos y, sin embargo, comportarse de manera muy
distinta para el stock: uno puede alternar arriba y abajo, mientras otro queda
corto de manera persistente.

La corrección multiplicativa de un segmento se estima con errores de A009 y
A010, sólo si sus sesgos tienen el mismo signo:
\begin{equation}
 f_s=\operatorname{clip}_{[0{,}8;\,1{,}2]}
 \left(\frac{\sum O_{s,t}}{\sum\widehat D_{s,t}}\right).
\end{equation}
Ese factor se prueba en A011. Se acepta para A012 y A013 únicamente si mejora
su error absoluto y el signo sigue siendo estable. El factor aceptado sigue
siendo el estimado \emph{antes} de A011; no se reajusta con el año que lo verifica.

\tabla{sesgo.tex}{Contraste de la corrección de sesgo. Los MAE corresponden a
A011; el factor se estimó con años anteriores.}
{tab:sesgo}{\archivo{resultados/demanda/tablas/diagnostico_sesgo_oof.csv}}

Se acepta una reducción de aproximadamente el 15,6\,\% en MD y un incremento
del 20\,\% en SSUV. En B, C y S, la corrección propuesta empeora A011 y se
descarta. MSUV no aporta un factor estable distinto de uno. La dispersión para
simular conserva los errores brutos de los ejercicios de selección, sin escoger
a posteriori sólo los años favorables; por ello tampoco certifica ausencia de sesgo.

\subsection{Qué nos dice A012 y qué no nos dice}
El error absoluto medio mensual es
\begin{equation}
 \mathrm{MAE}_s=\frac{1}{12}\sum_t|\widehat D_{s,t}-O_{s,t}|,
 \qquad
 \mathrm{sesgo}_s=\frac{1}{12}\sum_t(\widehat D_{s,t}-O_{s,t}).
\end{equation}
En el total comparable se suman los seis segmentos de cada mes \emph{antes}
de calcular el error. El resultado es MAE de \num{240.5} coches por mes, sesgo
de $-\num{190.2}$ y MAE relativo de 14,8\,\%.

\tabla{validacion.tex}{Evaluación retrospectiva A012, clientes observados.
Sesgo positivo: sobreprevisión. CSUV no forma parte del total.}
{tab:validacion}{\archivo{resultados/demanda/tablas/evaluacion_retrospectiva_a012.csv}}

El 14,8\,\% no debe compararse directamente con el 23,8\,\% de selección:
cambian los años y la agregación. Los errores de distintos segmentos pueden
compensarse en el total. Por ejemplo, el error relativo de C es 37,0\,\%, mucho
mayor que el del total agregado.

A012 ya se conocía durante el desarrollo inicial del proyecto. La reconstrucción
del protocolo elimina fugas concretas, pero no transforma ese año en una prueba
nunca vista. Se presenta como \textbf{evaluación retrospectiva}. Además se
vuelven a estimar previsiones desde cierres mensuales y se evalúan a hasta tres
meses; ese ejercicio responde al uso operativo de la revisión, no al plan anual
inicial.

\grafosanchos{validacion_a012_ocho_paneles.png}
{Clientes observados y previsión al cierre de A011. Las diferencias entre segmentos
explican por qué el MAE total no basta para fijar una única protección de stock.
El panel de CSUV muestra el lanzamiento, no una validación equivalente.}
{fig:validacion}{\archivo{demanda.ipynb}; evaluación retrospectiva A012}

\subsection{Cuota reciente, producto y el caso especial de CSUV}
La relación comercial acumula los últimos doce meses, excluyendo sucesos. Se usa
un cociente de sumas, no una media simple de porcentajes mensuales:
\begin{equation}
 q_s=\frac{\sum_{t\in\mathcal V_s}O_{s,t}}
           {\sum_{t\in\mathcal V_s}I_{s,t}},
\end{equation}
donde $\mathcal V_s$ identifica los meses seleccionados; se evita dar el mismo
peso a meses con mercados muy distintos.

\paso{Producto.}
Se comparan ventanas completas y observables de seis meses antes y después de
renovaciones, sin sucesos, roturas ni episodios solapados. El ratio utiliza clientes
netos sobre mercado. Al no aislar causalidad, el efecto central adicional es cero.
El optimista aplica una sensibilidad empírica tras retirar la exposición hipotética
ya contenida en el ancla, evitando duplicidades.

Se supone un calendario de enero y seis meses de efecto para B, SSUV, CSUV y
MSUV. Con efecto central cero, retrasar MSUV hasta abril no cambia por sí solo
la cantidad central; la pérdida de cuota por retraso se estudia en una prueba
específica de sensibilidad. \emph{No identificado} no significa
\emph{económicamente irrelevante}.

\paso{CSUV.}
Se utilizan septiembre--diciembre, excluyendo agosto parcial y ajustando los picos;
se contrasta septiembre--octubre. El mercado permite proyectar todo A013 sin
repetir ceros anteriores al lanzamiento. Su dispersión es un proxy de SSUV/MSUV
ampliado por 1,5, con sensibilidad 1--2; no un error propio validado.

\subsection{La salida de esta fase: doce meses, siete segmentos}
El central combina mercado previsto, relación comercial y factores de sesgo
aceptados. Después se descuenta el adelanto A012 una sola vez, repartido en los
tres primeros meses. Los escenarios de demanda cambian macroeconomía, ayudas,
producto, turismo y duración de la devolución.

\tabla{forecast.tex}{Demanda anual prevista antes del efecto residual de combustible.
Son previsiones condicionadas, no ventas simuladas ni medias de réplicas.}
{tab:forecast}{\archivo{resultados/demanda/tablas/prevision_anual_por_escenario.csv}}

El central estima unos \num{19077} clientes. El pesimista y optimista de esta
fase estiman aproximadamente \num{16210} y \num{27700}. En
\archivo{escenarios.ipynb} se añade el efecto residual del combustible antes
del descuento absoluto del adelanto. Por eso no deben confundirse estas cifras
intermedias con la demanda final de los escenarios ni con las ventas alcanzadas.

\FloatBarrier
\grafosanchos[\textbf{Del pronóstico a la logística.} Esta fase entrega una
trayectoria central, alternativas y errores, no un pedido automático. El simulador
descuenta stock y compromisos, añade protección y calcula las consecuencias.
La incertidumbre de CSUV y los sesgos residuales siguen formando parte de la decisión.]
{prevision_a013_ocho_paneles.png}
{Previsiones condicionadas A013 y banda P10--P90 del central. El optimista puede
tener un enero más bajo al concentrar la devolución del adelanto en ese mes,
aunque termine el año con mayor demanda. Las escalas varían por segmento.}
{fig:forecast}{\archivo{demanda.ipynb}; 1.000 trayectorias ilustrativas, semilla 42}

\section{Tercera fase: convertir demanda en pedidos, stock y costes}\label{sec:simulacion}

\begin{idea}{Entrada, trabajo y salida}
\textbf{Entran} previsión central, stock inicial, llegadas comprometidas, errores
de previsión y catálogo económico. \textbf{Se simula} mes a mes el movimiento
de cohortes, las ventas, su antigüedad y las decisiones a tres meses.
\textbf{Salen} cantidades físicamente consistentes y una cuenta económica.
Primero se comprueba una trayectoria sin ruido; después se repite el mecanismo.
\end{idea}

\subsection{Pedido, fabricación asignada y llegada: tres fechas}
El caso aproxima el plazo logístico por tres meses. Si $Q_{s,t}$ es el pedido
decidido al cierre del mes $t$ y $L_{s,t+3}$ la llegada correspondiente,
\begin{equation}\label{eq:plazo}
 L_{s,t+3}=Q_{s,t}
\end{equation}
cuando no existe una interrupción. La fabricación se aproxima por la asignación
ejecutada del pedido; no se representa una planta con su propia capacidad y
secuencia diaria.

Al cerrar diciembre de A012, enero y febrero ya están comprometidos. Marzo es
la primera llegada modificable. Una revisión al cerrar enero modifica abril,
no febrero ni marzo. La distinción evita instrucciones como «pedir en septiembre»
cuando en realidad se necesita \emph{recibir} en septiembre.

\begin{figure}[htbp]\centering
\begin{tikzpicture}[x=2.15cm,y=1cm,font=\small,>=Stealth]
\foreach \x/\mes in {0/dic A012,1/ene A013,2/feb A013,3/mar A013,4/abr A013,5/may A013}{
  \node at (\x,0) {\mes};\draw[gris!50] (\x,-.15)--(\x,-.3);}
\draw[gris!50] (0,-.3)--(5,-.3);
\draw[->,azul,thick] (0,-.65)--node[below]{decisión de diciembre: llegada en marzo}(3,-.65);
\draw[->,azul,thick] (1,-1.5)--node[below]{revisión de enero: llegada en abril}(4,-1.5);
\node[draw=gris!40,fill=fondo,align=center,text width=3.4cm] at (1.5,.9)
  {Enero y febrero\\ya vienen en camino};
\end{tikzpicture}
\caption{El plazo de decisión. El stock disponible y los coches en camino no son
la misma cantidad. Los pedidos de octubre--diciembre de A013 llegarán en A014.}
\label{fig:plazo}\fuente{regla de la sección 7 del enunciado y API logística revisada}
\end{figure}

\subsection{Cómo se fija el estado inicial}
El stock al final de A012 es observado: \num{2877} coches. Las llegadas de enero
y febrero no se conocen directamente y se estiman con la media mensual de
\emph{wholesale} de los últimos seis meses, desde que existe actividad en el
segmento. Se redondean a enteros y se suponen iguales para ambos meses.

\tabla{inicial.tex}{Estado al empezar A013. El stock es observado; el tránsito
es estimado. Las coberturas usan la previsión mensual media A013 como denominador.}
{tab:inicial}{\archivo{resultados/simulacion/tablas/estado_inicial_a013.csv}; total calculado con sumas}

El total supuesto en camino para enero--febrero es \num{4220} coches. Por eso la
cobertura del inventario disponible es de 1,8 meses, pero al incorporar tránsito
sube a 4,5. No se puede usar la segunda cifra para afirmar que todos esos coches
pueden venderse hoy. C ilustra el problema contrario: parte ya de un inventario
alto, que limita la necesidad de pedidos aunque siga existiendo demanda.

Las edades iniciales se representan con lotes de 15, 45 y 75 días; 105, 135 y
165 días; y 210 días para el tramo superior a 180. Se conservan unidades enteras.
Las sensibilidades desplazan las edades 15 días y comparan referencias de tránsito
de tres o seis meses, o cantidades un 25\,\% menores o mayores.

\subsection{El balance físico de un mes}
Sin transfers, con demanda realizada $D_{s,t}$ y stock anterior $S_{s,t-1}$:
\begin{align}
 V^{\mathrm{red}}_{s,t}&=\min(D_{s,t},S_{s,t-1}+L_{s,t}),\\
 S_{s,t}&=S_{s,t-1}+L_{s,t}-V^{\mathrm{red}}_{s,t}.
\end{align}
Las llegadas se consideran disponibles al inicio del mes y las ventas se imputan
tras 15 días. Es una convención mensual para valorar edad y coste, no una
simulación de cada comprador durante el día.

\paso{Ejemplo pequeño.}
Hay 20 coches y llegan 50. Si 80 clientes quieren comprar, se venden 70 de la
red y quedan diez sin atender. Con cinco transfers directos, las ventas totales
son 75 y se pierden cinco. El stock final de la red es cero en ambos casos:
los transfers no deben restarse otra vez de un stock al que no se incorporaron.

En la alternativa de transfer se utiliza una fracción $a$ del déficit:
\begin{align}
 T_{s,t}&=\left\lfloor a\,\bigl(D_{s,t}-V^{\mathrm{red}}_{s,t}\bigr)\right\rfloor,\\
 V_{s,t}&=V^{\mathrm{red}}_{s,t}+T_{s,t},\qquad
 P_{s,t}=D_{s,t}-V_{s,t},
\end{align}
con $P$ igual a ventas perdidas. El caso base usa $a=0$ porque no hay
disponibilidad garantizada. Se contrasta $a=0{,}25$ y $a=1$ como referencia
favorable; no se asume que los países europeos tengan siempre stock ilimitado.

\subsection{Por qué FIFO cambia el resultado económico}
El balance total no informa de la edad de los coches que se venden. FIFO
significa que se atiende la demanda con las cohortes más antiguas primero.
Dos políticas con el mismo stock final pueden generar gastos comerciales
distintos si conservan edades diferentes.

En cada mes se añade la cohorte nueva, se venden primero las antiguas, se
imputa el VM de cada edad y se envejece lo que queda. Las cohortes iniciales
mantienen sus días anteriores a enero: un coche con 170 días no recibe otros
180 días subvencionados al cambiar de año.

\subsection{Coste de almacenamiento e intereses}
Para un coche de edad inicial $a$ que permanece otros $d$ días, el coste adicional
de la marca es
\begin{equation}\label{eq:coste-stock}
 C^{\mathrm{marca}}(a,d)
 =2\left[\min(a+d,180)-\min(a,180)\right]\unidad{\euro}.
\end{equation}
Así, un coche recién llegado que se vende en el mismo mes genera
$2\times15=30$ euros. Un coche de 170 días que permanece 30 días adicionales
genera sólo $2\times10=20$ euros. El máximo de toda su vida facturada es
360 euros; no se extrapola el coste durante años.

\begin{figure}[htbp]\centering
\begin{tikzpicture}
\begin{axis}[width=.88\linewidth,height=6cm,xmin=0,xmax=300,ymin=0,ymax=410,
  xlabel={Edad total del coche (días)},ylabel={Coste acumulado de la marca (\euro)},
  xtick={0,60,90,120,180,240,300},ytick={0,120,240,360},legend pos=south east]
\addplot[azul,thick,domain=0:300,samples=101]{2*min(x,180)};
\addlegendentry{Regla del caso}
\addplot[gris,dashed] coordinates {(180,0) (180,360)};
\end{axis}\end{tikzpicture}
\caption{El coste de la marca se detiene a los 180 días. El gasto comercial por
edad y los intereses del concesionario son conceptos adicionales.}
\label{fig:coste-stock}\fuente{enunciado, sección 2.4; ecuación~\eqref{eq:coste-stock}}
\end{figure}

Los intereses del concesionario se calculan por tramos: tarifa diaria del primer
tramo hasta el día 90, del segundo entre 91 y 180, y ningún interés adicional
contado por el caso después. Si $c_1,c_2$ son las tarifas diarias del catálogo,
la función acumulada es
\begin{equation}
 J(a)=c_1\min(a,90)+c_2\min\!\left(\max(a-90,0),90\right).
\end{equation}
El coste de un periodo es $J(a+d)-J(a)$. Por ejemplo, pasar de 85 a 100 días
supone cinco días al primer tipo y diez al segundo, no quince días a ambos.
Estos intereses se muestran separados del resultado de la marca. El vencimiento
del crédito a los 180 días es un riesgo de liquidez para el concesionario.

\subsection{Cuánto pedir y cómo dimensionar el colchón}
Al cerrar $t$, se proyecta el stock disponible antes de la primera llegada
modificable, usando las llegadas ya comprometidas y la previsión de $t+1$ y
$t+2$. La implementación no permite stock proyectado negativo: una demanda que
no se atiende no se convierte en un inventario negativo a reponer, porque el
caso base no acumula pedidos pendientes de clientes.

La regla de pedido es
\begin{equation}\label{eq:pedido}
 Q_{s,t}=\left\lceil\max\!\left(0,
 \widehat D_{s,t+3}+SS_{s,t}-\widehat S_{s,t+2}\right)\right\rceil.
\end{equation}
El redondeo hacia arriba expresa coches enteros. $SS$ es el stock de seguridad;
no es una reserva idéntica para todos los segmentos.

La fórmula docente $z\sigma\sqrt{3}$ supone errores mensuales independientes.
La revisión utiliza errores acumulados observados en ventanas de tres meses:
\begin{equation}
 \sigma_{3,s}^{\mathrm{rel}}
 =\operatorname{sd}\left(
 \frac{\sum_{u=t+1}^{t+3}(\widehat D_{s,u}-O_{s,u})}
      {\sum_{u=t+1}^{t+3}\widehat D_{s,u}}\right),\qquad
 SS_{s,t}=z_s\,\sigma_{3,s}^{\mathrm{rel}}
                \sum_{u=t+1}^{t+3}\widehat D_{s,u}.
\end{equation}
Las ventanas no se unen entre años de validación distintos. El error ya está
acumulado durante el plazo; multiplicar otra vez por $\sqrt3$ duplicaría esa
adaptación. Se cubren los meses \emph{hasta la llegada}, no los tres posteriores
a recibirla.

\paso{Comprobación del ejemplo del enunciado.}
Con 536 coches iniciales, 700 de llegada en enero y otros 700 en febrero,
y previsiones 550 y 600, el stock proyectado es $536+700+700-550-600=786$.
Con $\sigma=100$ y $z=1{,}28$, el colchón docente es aproximadamente 222 coches.
Para una previsión de marzo de 850, resulta
\[
 Q=\left\lceil850+1{,}28\times100\sqrt3-786\right\rceil=286.
\]
El notebook comprueba este valor ejercitando el planificador, además de verificar
balances y fronteras de coste.

\subsection{Revisar sin conocer la demanda que no se atendió}
La revisión utiliza el cociente entre las ventas y la previsión acumuladas en
los tres últimos meses, $r$. La corrección amortiguada es
\begin{equation}\label{eq:revision}
 f=\max\bigl(0,\,1+0{,}5(r-1)\bigr).
\end{equation}
Si las ventas superan un 20\,\% la previsión, el factor es 1,10: se incorpora
la mitad de la desviación para evitar reaccionar plenamente a un mes ruidoso.

Si la red agotó su disponibilidad en la ventana, las ventas están censuradas
por falta de coches. En ese caso el factor no se reduce por debajo de uno.
El planificador observa ventas y disponibilidad, no utiliza mágicamente la
demanda latente simulada para corregirse. El factor se aplica a las previsiones
del plazo y sólo cambia el pedido que llegará en $t+3$.

\subsection{Primero comprobar un año sin ruido}
La referencia determinista utiliza la previsión central redondeada al entero
más próximo, con $z=1{,}28$. No es todavía la política recomendada ni una media
de mil simulaciones.

\subsubsection*{Un recorrido real: SSUV en el primer trimestre}
\begin{table}[H]\centering
\caption{Primer trimestre SSUV de la referencia determinista, coches.}
\begin{tabular}{@{}lrrrrrr@{}}\toprule
Mes & \makecell{Stock\\inicial} & Llegadas & Demanda & Ventas & Pérdidas & \makecell{Stock\\final} \\\midrule
Enero & 0 & 390 & 369 & 369 & 0 & 21 \\
Febrero & 21 & 390 & 500 & 411 & 89 & 0 \\
Marzo & 0 & 1.420 & 591 & 591 & 0 & 829 \\\bottomrule
\end{tabular}
\fuente{\archivo{resultados/simulacion/tablas/simulacion_central_mensual_a013.csv}}
\end{table}

En enero, $0+390-369=21$ coches quedan en stock. En febrero sólo hay
$21+390=411$ para una demanda de 500, de modo que se pierden 89 compras.
La llegada de marzo todavía no está disponible para esos clientes. Al llegar
los \num{1420} coches de marzo, la demanda es 591 y quedan 829 de protección;
las compras perdidas de febrero no se convierten en pedidos pendientes.

El coste de enero también puede comprobarse a mano: 369 coches vendidos con
15 días y 21 que permanecen 31 días generan
$369\times15\times2+21\times31\times2=\num{12372}$ euros. La convención mensual
es la misma que se usa después en todas las réplicas aleatorias.

\tabla{determinista.tex}{Referencia determinista A013: demanda central sin ruido,
colchón uniforme $z=1{,}28$. Resultado tras costes modelizados.}
{tab:determinista}{\archivo{resultados/simulacion/tablas/resumen_central_a013.csv}}

Atiende \num{18991} coches y deja \num{89} sin atender, todos en SSUV.
La contribución es 64,18 M\euro{} y el stock final, \num{2319}. Tener prevista
la demanda no elimina las restricciones iniciales: en SSUV se parte de stock
cero y las dos primeras llegadas ya están comprometidas.

\FloatBarrier
\grafosanchos[\textbf{Antes de añadir azar.} Se verifican balances y unidades
enteras, FIFO y VM por edad, los 15 días de una venta nueva, el límite de 180 días,
intereses por tramos, transfers, ausencia de anticipación y el pedido de diciembre
que llega en marzo A014. Repetir una simulación no corrige un mecanismo mal definido.]
{stock_central_a013_a014.png}
{Referencia determinista y prolongación hasta junio de A014. El punto gris es
stock inicial observado; el trazo azul, stock simulado. La línea vertical separa
A013 de la sensibilidad de continuidad, que no es una previsión validada.}
{fig:simulacion}{\archivo{simulacion.ipynb}; política base, escalas por segmento}

\section{Cuarta fase: escenarios, incertidumbre y elección de política}\label{sec:escenarios}

\begin{idea}{Dos preguntas que requieren respuestas distintas}
\textbf{Escenario:} ¿qué cambia si la economía, la normativa o el suministro son
distintos? \textbf{Incertidumbre dentro del escenario:} aun bajo esos supuestos,
¿cuánto pueden variar los meses? Las sensibilidades examinan además si una
hipótesis poco identificada cambia la decisión. No son tres nombres para la misma cosa.
\end{idea}

\subsection{De una historia a parámetros concretos}
El central mantiene las condiciones económicas del cierre de A012. Los extremos
utilizan un modelo macroeconómico auxiliar para obtener un único cociente entre
mercado bajo el escenario y mercado auxiliar central. Ese cociente se aplica a
la previsión elegida: no se sustituye el central por otro modelo ni se suma una
penalización económica genérica adicional.

\begin{table}[htbp]\centering
\caption{Principales supuestos A013. Las variaciones de indicadores se miden
respecto al cierre de A012; no son variaciones de ventas.}\label{tab:escenarios-param}
\small
\begin{tabularx}{\linewidth}{@{}p{3.6cm}YYY@{}}\toprule
Factor & Pesimista & Central & Optimista \\\midrule
Paro & +2 puntos porcentuales & Sin cambio & $-1$ punto porcentual \\
Confianza & $-16{,}26$ puntos de índice & Sin cambio & +19,14 puntos de índice \\
Tipo de crédito & +0,546 puntos porcentuales & Sin cambio & $-0{,}546$ puntos porcentuales \\
Devolución del adelanto & 6 meses & 3 meses & 1 mes \\
Ayuda pública & No & No & Abril--diciembre \\
Producto & Central neutro; retraso en sensibilidad & Efecto adicional cero & +16,8\,\% hipotético durante 6 meses, con normalización previa \\
Turismo & Cambio bajo de viajeros & Sin cambio & Cambio alto de viajeros \\
Combustible residual & $-3\,\%$ de demanda & Cero & $+3\,\%$ de demanda \\
CO$_2$ & Prima del catálogo por venta & Sin prima & Sin prima \\
Nueva importación & +500 \euro/coche por arancel & Sin recargo & $-250$ \euro/coche por divisa \\
Suministro MSUV & Pérdida de producción en marzo; llegada junio & Sin interrupción & Sin interrupción \\\bottomrule
\end{tabularx}
\fuente{\archivo{resultados/escenarios/tablas/parametros_escenarios_a013.csv}}
\end{table}

Los choques de confianza y viajeros se apoyan en percentiles de cambios históricos
a nueve meses; el del crédito, en la mediana absoluta de esos cambios. El paro
sigue la propuesta del enunciado. Turismo afecta al canal de alquiler durante
marzo--junio, ponderado por su peso en los clientes de Kaido, no por el peso del
alquiler en todo el mercado.

El combustible, el arancel y el cambio no tienen una base suficiente para estimar
sus tamaños causalmente. Sus importes son \textbf{pruebas de estrés explícitas},
contrastadas con valores cero y dobles. El coste arancelario no se calcula como
un porcentaje arbitrario del PVP. La divisa sólo cambia el coste de las nuevas
importaciones; no revaloriza gratuitamente los coches iniciales.

\subsection{En qué orden se aplican los efectos}
La composición importa. Primero se aplican efectos sobre mercado y posición,
turismo y combustible. Después se descuenta el adelanto expresado en coches:
\begin{equation}\label{eq:demanda-escenario}
 \widehat D^{(e)}_{s,t}
 =\max\!\left(0,\ \widehat D^{(e),\mathrm{previa}}_{s,t}(1+c_e)
                       -\frac{A_s}{m_e}\,\mathbf1_{t\leq m_e}\right).
\end{equation}
$c_e$ es el efecto residual del combustible; $A_s$, el adelanto A012; y $m_e$,
el número de meses de devolución. La misma cantidad central $A_s$ se solicita
en los tres escenarios. Multiplicar después por $1+c_e$ alteraría también el
adelanto absoluto, que no debe convertirse en una segunda elasticidad.

En el optimista, concentrar la devolución en enero supera la demanda prevista
de algunos segmentos. El límite a cero deja \num{235.3} coches del adelanto sin
descontar en la integración final. Se declara esa parte no aplicada; no se crea
demanda negativa ni se traslada silenciosamente a febrero. Es una limitación de
la combinación de hipótesis, no un error de redondeo.

\subsection{Generar años posibles sin perder la dependencia entre meses}
Una réplica produce 18 meses de demanda realizada: A013 y seis meses de
continuidad. La realización se genera mediante
\begin{equation}\label{eq:ruido}
 D^{(r,e)}_{s,t}=\operatorname{redondear}
 \left[\widehat D^{(e)}_{s,t}
       \exp\!\left(\sigma_{\log,s}Z^{(r)}_{s,t}
                         -\frac{\sigma_{\log,s}^{2}}{2}\right)\right].
\end{equation}
El factor exponencial es positivo. Si $Z$ es normal de media cero y varianza uno,
restar $\sigma_{\log,s}^{2}/2$ hace que el factor tenga media uno antes del
redondeo. No significa que cada réplica coincida con la previsión ni que su
mediana sea igual a su media.

El ruido contiene una parte común y una específica:
\begin{align}
 Z_{s,t}&=\sqrt\rho\,U_t+\sqrt{1-\rho}\,W_{s,t},\\
 U_t&=\phi U_{t-1}+\sqrt{1-\phi^2}\,\eta_t,
\end{align}
con una ecuación análoga para $W$. Las innovaciones $\eta$ son normales
independientes; $\phi$ introduce persistencia temporal y $\rho$ co-movimiento
entre segmentos. En la ejecución, $\phi\simeq0{,}210$ y $\rho\simeq0{,}190$.
No se asume que un mes alto se compense necesariamente con uno bajo.

\tabla{incertidumbre.tex}{Dispersión de errores por segmento. Son parámetros
adimensionales; $\sigma_{3,\mathrm{rel}}$ corresponde a errores acumulados en tres meses.}
{tab:incertidumbre}{\archivo{resultados/escenarios/tablas/parametros_incertidumbre_a013.csv}}

Los parámetros proceden de los errores de selección en segmentos maduros,
separando años al calcular ventanas y correlaciones. El logaritmo de error usa
un suavizado de 0,5 coche para tratar ceros. CSUV recibe el proxy ampliado. El
valor muy elevado de la dispersión relativa de S revela una serie pequeña e
inestable: no debe interpretarse como una precisión demostrada para ese nicho.

\subsection{Qué ocurre dentro de cada réplica}
La secuencia se repite sin cambiar las reglas de contabilidad:
\begin{enumerate}
  \item Se parte del mismo stock y del mismo plan inicial para los tres escenarios.
  \item Llegan los compromisos del mes, respetando las interrupciones conocidas.
  \item Se atiende la demanda con FIFO y, sólo en la alternativa correspondiente,
  con transfers disponibles.
  \item Se anotan ventas, faltantes, edades, stock y costes.
  \item Se observan ventas y disponibilidad, se revisa la previsión y se decide
  únicamente la llegada de dentro de tres meses.
  \item Se continúa hasta junio de A014 para examinar los efectos de cierre.
\end{enumerate}

La falta de una pieza en fábrica en marzo afecta a la llegada de junio de MSUV.
No borra coches que ya estaban viajando. La producción perdida no se recupera
automáticamente después; hacerlo requeriría otra hipótesis de capacidad.

Todos los escenarios utilizan las mismas innovaciones aleatorias al compararse.
Esto no impone probabilidades a los escenarios: reduce ruido al comparar políticas.
La semilla 42 permite reproducir las trayectorias.

\subsection{Cómo leer media, mediana y P10--P90}
\begin{description}[style=nextline]
\item[Media.] Promedio de las réplicas; es la magnitud usada para la contribución
esperada bajo un escenario.
\item[Mediana.] Mitad de las réplicas por encima y mitad por debajo.
\item[P10 y P90.] Aproximadamente el 10\,\% de las réplicas queda por debajo del
primero y el 10\,\% por encima del segundo. El intervalo intermedio contiene
aproximadamente el 80\,\% de las realizaciones de la simulación.
\end{description}
No son un intervalo de confianza para una verdad desconocida ni una probabilidad
de que se produzca el escenario. Para el total se suman los segmentos
\emph{dentro de cada réplica} y después se calculan sus percentiles. La suma de
medianas de segmentos no es, en general, la mediana del total.

\subsection{Elegir una política: margen sujeto a servicio}
El servicio de un segmento en una réplica es la proporción de su demanda anual
que se ha podido atender:
\begin{equation}\label{eq:servicio}
 F_s^{(r)}=\frac{\sum_{t\in A013}V_{s,t}^{(r)}}
                 {\sum_{t\in A013}D_{s,t}^{(r)}}.
\end{equation}
La restricción acordada es que la \textbf{media de $F_s$ en el central} alcance
el 95\,\% en los segmentos marginalmente viables. No es un 95\,\% de
probabilidad de pasar el año sin roturas ni un requisito en cada escenario.

Se prueban seis niveles de protección:
$z\in\{0;\,0{,}84;\,1{,}28;\,1{,}65;\,2{,}33;\,3\}$.
El piloto usa 200 réplicas. Como no se impone una capacidad compartida de fábrica
no cuantificada en el caso, cada segmento puede escoger su nivel.

Entre políticas factibles se busca mayor contribución media central. Cuando una
opción con menor colchón queda a no más de dos errores estándar de Monte Carlo
de la mejor, se prefiere la menor protección. El error estándar se calcula sobre
las diferencias emparejadas:
\begin{equation}
 \mathrm{EE}(\overline{\Delta R})
 =\frac{\operatorname{sd}(R^{(r)}_{\mathrm{mejor}}-R^{(r)}_z)}{\sqrt N}.
\end{equation}
Esta tolerancia depende de la precisión computacional de la comparación, no
de una cantidad fija idéntica para todos los segmentos. Tampoco mide toda la
incertidumbre económica del modelo.

La política elegida se comprueba con 1.000 réplicas. El código permite aumentar
el colchón si esa verificación descubre un incumplimiento; en la ejecución
presentada no fue necesario cambiar los valores elegidos por el piloto.

\tabla{politica.tex}{Política seleccionada y servicio anual central. La
probabilidad indicada es la frecuencia simulada de $F_s<95\,\%$, condicionada al central.}
{tab:politica}{\archivo{resultados/escenarios/tablas/politica_recomendada_a013.csv}}

SSUV alcanza un servicio medio del 95,4\,\%, pero en el 39,1\,\% de las
réplicas no llega al 95\,\%. No hay contradicción: la restricción está formulada
sobre la media. Su P10 del 89,3\,\% expresa un riesgo que la media oculta.

\begin{figure}[htbp]\centering
\begin{tikzpicture}
\begin{axis}[width=.94\linewidth,height=6.7cm,xmin=-.4,xmax=6.4,ymin=80,ymax=101,
  xtick={0,1,2,3,4,5,6},xticklabels={B,C,MD,SSUV,CSUV,MSUV,S},
  ylabel={Demanda anual atendida (\%)},xlabel={Segmento},
  legend style={at={(.5,-.24)},anchor=north,legend columns=3}]
\addplot[azul,only marks,mark=*,mark size=2.4pt]
  table[col sep=comma,x expr=\coordindex,y expr=100*\thisrow{servicio_central}]
  {datos/politica_recomendada_a013.csv};\addlegendentry{Media central}
\addplot[gris,only marks,mark=square*,mark size=2.1pt]
  table[col sep=comma,x expr=\coordindex,y expr=100*\thisrow{p10_servicio}]
  {datos/politica_recomendada_a013.csv};\addlegendentry{P10 central}
\addplot[azul,dashed,domain=-.4:6.4]{95};\addlegendentry{Objetivo en media}
\end{axis}\end{tikzpicture}
\caption{Cumplir el objetivo medio no elimina años desfavorables.
El eje vertical está ampliado al intervalo 80--101\,\%; son puntos, no barras
de cantidades con origen recortado.}
\label{fig:servicio}\fuente{política recomendada, 1.000 réplicas centrales}
\end{figure}

\begin{idea}{Alcance de la elección}
Se recomienda la mejor opción entre las políticas comparadas bajo el objetivo
declarado. No se demuestra un óptimo global ni una política óptima para cualquier
probabilidad de escenarios. La malla, la ausencia de capacidad compartida y la
forma de los errores son parte del experimento.
\end{idea}
\FloatBarrier

\section{Leer los resultados y convertirlos en una recomendación}\label{sec:lectura-final}

\subsection{Antes de sumar euros: qué cuenta cada concepto}
El catálogo separa beneficio bruto de la marca, VM y CVP. En la simulación,
el VM depende de la edad del coche vendido:
\begin{align}
 \mathrm{CVP}&=\mathrm{beneficio\ bruto}-\mathrm{VM},\\
 R&=\mathrm{CVP}-C_{\mathrm{stock}}-C_{\mathrm{transfer}}
       -C_{\mathrm{CO_2}}-C_{\mathrm{arancel}}-C_{\mathrm{cambio}}.
\end{align}
$R$ es contribución después de los costes modelizados, no beneficio neto contable
de una empresa con todos sus gastos generales. Un coste de cambio negativo
representa un ahorro y, al restarlo, aumenta el resultado.

No se vuelve a restar VM después de calcular CVP. Tampoco se descuentan otra vez
las ventas perdidas: sus ingresos ya no forman parte de las ventas realizadas.
El coste de oportunidad se informa aparte. Los aranceles y el cambio adicional
se imputan a las ventas de cohortes de nueva importación; el coste incremental
que permanece en el stock final se registra, pero no se convierte en ingreso.

\tabla{unidad.tex}{Economía de una unidad joven por segmento, antes de stock y
otros recargos. Todos los importes son euros por coche; el coste de transfer
es el 20\,\% del beneficio bruto, no del CVP.}
{tab:unidad}{\archivo{resultados/escenarios/tablas/economia_unitaria_segmento.csv}}

Un MSUV joven deja \num{5887} euros de CVP antes de costes logísticos; mantener
un coche nuevo 15 días cuesta 30 euros. No es una comparación que pueda
extenderse sin límite: el gasto de stock de la marca se detiene en 360 euros y
el VM de una unidad vieja puede ser diferente.

En B el CVP joven es sólo 125 euros antes de stock, frente a una prima CO$_2$
de \num{2555}. En C, el CVP joven de \num{1367} pasa a $-585$ euros al cargar
CO$_2$, antes de arancel y stock. Una venta adicional no siempre mejora la
contribución. Por eso, una vez confirmados los costes del escenario, se frenan
\emph{nuevas reposiciones} marginalmente negativas. No se deja de vender
automáticamente el inventario ya comprometido.

\subsection{Resultado agregado de la política recomendada}
La tabla~\ref{tab:resultado} utiliza medias de 1.000 réplicas. En particular,
\emph{pedidos A013} cuenta decisiones de ese año, mientras \emph{llegadas}
incluye coches comprometidos en A012. No tienen por qué coincidir.

\tabla{resultados.tex}{Política recomendada: resultados medios A013, total de la
marca. Las coberturas usan demanda mensual media de cada réplica.}
{tab:resultado}{\archivo{resultados/escenarios/tablas/tabla5_plan_recomendado_a013.csv}}

\paso{Central.}
Se simulan unas \num{18539} ventas, \num{478} compras perdidas y un servicio
medio del 97,6\,\%. La contribución media es 62,88 M\euro{}. El stock final
medio es \num{2221} coches. La demanda media simulada de \num{19017} no tiene
por qué igualar exactamente la previsión puntual de \num{19077}: procede de
una muestra finita de trayectorias y del redondeo a unidades.

\paso{Pesimista.}
La contribución cae a 3,09 M\euro{} y el servicio al 77,8\,\%. No se debe sólo
a una demanda menor. CO$_2$ y el recargo de importación hacen negativas nuevas
reposiciones de B, C, SSUV y S bajo la regla marginal, y se incorpora la
interrupción de suministro. El deterioro de servicio es una consecuencia
importante de la política económica, no una cifra que pueda ocultarse con el
cumplimiento del central.

\paso{Optimista.}
Se simulan unas \num{24952} ventas y 91,37 M\euro{} de contribución media.
Sin embargo, se pierden unas \num{3493} compras: la demanda mejora más deprisa
que la capacidad de modificar llegadas. El ahorro de divisa, registrado como
coste negativo, explica que la contribución final pueda superar el CVP antes
de esos ajustes. No se han restado los costes dos veces ni creado ventas extra.

\tabla{riesgo.tex}{Dispersión del resultado total y del servicio anual.
P10--P90 describe variación dentro del escenario, no probabilidades de escenarios.}
{tab:riesgo}{\archivo{resultados/escenarios/tablas/tabla5_plan_recomendado_a013.csv}}

El central tiene P10--P90 de contribución de 56,45--69,47 M\euro{}. Es un rango
condicionado a los supuestos y a la distribución de errores elegida. El error
estándar de su media de Monte Carlo es aproximadamente 0,159 M\euro{}, una
medida de precisión de simulación que no equivale a ese rango de resultados.
No observar pérdidas en las 1.000 réplicas finales no demuestra que su
probabilidad real sea cero.

\grafica{plan_recomendado_ventas_stock_a013.png}
{Ventas y stock del plan recomendado: mediana y P10--P90 por escenario.
Las bandas reflejan variación mensual de una misma política inicial y sus
revisiones. El stock representado no incluye coches en camino.}
{fig:resultado-mensual}{\archivo{escenarios.ipynb}; 1.000 réplicas por escenario}

\subsection{No todos los segmentos desempeñan el mismo papel}
\tabla{segmentos.tex}{Resultado central por segmento: medias de demanda, ventas,
contribución, stock final y servicio. Las sumas pueden diferir por redondeo.}
{tab:segmentos}{\archivo{resultados/escenarios/tablas/tabla5_plan_recomendado_a013.csv}}

MSUV aporta aproximadamente 31,91 M\euro{}, cerca de la mitad del resultado
central. SSUV y CSUV aportan 14,86 y 11,77 M\euro{}, respectivamente. La
disponibilidad de esos segmentos tiene más impacto absoluto que una mejora
pequeña en un segmento de bajo volumen.

B apenas aporta 0,02 M\euro{} en el central: mantener su objetivo de servicio
no convierte su economía en robusta. C conserva un inventario inicial elevado;
una cobertura alta no justifica encargar más coches. S combina bajo volumen,
alta dispersión relativa y stock medio elevado frente a la demanda. Su cobertura
debe leerse como señal de exposición, no como garantía de calidad de previsión.

\subsection{Aislar el CO\texorpdfstring{$_2$}{2} sin confundir escenarios}
Para medir únicamente la prima, se toman las mismas ventas y flujos centrales
y se resta la prima del catálogo a cada venta. No se reoptimiza la política
en esta comparación.

\tabla{co2.tex}{Prima CO$_2$ aislada sobre los mismos flujos centrales.
Importes en M\euro{}; efecto negativo significa reducción del resultado.}
{tab:co2}{\archivo{resultados/escenarios/tablas/efecto_co2_aislado_a013.csv}}

La reducción media aislada es 55,01 M\euro{}, dejando 7,87 M\euro{} sobre esos
mismos flujos. El pesimista completo da 3,09 M\euro{} porque también cambia
el volumen servido, los pedidos, los aranceles y el suministro. No sería
correcto atribuir toda la diferencia central--pesimista a CO$_2$.

\subsection{Qué significa equivocarse por arriba o por abajo}
\begin{description}[style=nextline]
\item[Pasarse.] El exceso inmoviliza stock, consume días de coste y puede exigir
más VM al venderse viejo. La comparación debe considerar los costes posteriores
del inventario final, no sólo diciembre.
\item[Quedarse corto.] Se pierde la contribución de una venta potencial o se
incurre en el coste de conseguir un transfer. El margen no obtenido debe
descontar también CO$_2$, arancel y otros costes que no se pagan al no vender.
Puede resultar negativo en una venta marginal inviable.
\end{description}
El simulador no calcula un presupuesto de publicidad óptimo ni una respuesta de
demanda a cada euro de descuento. «Empujar» un segmento significa aquí priorizar
abastecimiento o revisar su política comercial, no asegurar que cualquier campaña
adicional sea rentable.

\subsection{La sensibilidad: comprobar qué hipótesis sostiene la decisión}
Las pruebas cambian un supuesto cada vez y utilizan las mismas 200 réplicas del
piloto. Por tanto, no hay que comparar mecánicamente sus importes con la media
final de 1.000 réplicas. La referencia central de la política final con esas
200 réplicas es 63,41 M\euro{}, frente a 62,88 M\euro{} con 1.000.

\begin{table}[htbp]\centering
\caption{Ejemplos de sensibilidad, 200 réplicas comunes.
Los importes son medias condicionales A013 y no se comparan con otra muestra.}
\label{tab:sensibilidad}
\begin{tabularx}{\linewidth}{@{}Yrr@{}}\toprule
Prueba & \makecell{Resultado\\(M\euro)} & \makecell{Servicio\\medio (\%)} \\\midrule
Central: referencia de las mismas 200 réplicas & 63,41 & 97,7 \\
Central: transfer del 25\,\% del déficit & 63,63 & 98,3 \\
Central: transfer suficiente, cota favorable & 64,29 & 100,0 \\
Pesimista: sin arancel adicional & 7,78 & 78,1 \\
Pesimista: arancel de 1.000 \euro/coche & $-1{,}50$ & 78,0 \\
Pesimista: no frenar reposición negativa & $-0{,}55$ & 98,6 \\
Central: demanda CSUV un 25\,\% menor & 60,34 & 97,9 \\
Central: demanda CSUV un 25\,\% mayor & 66,32 & 97,3 \\\bottomrule
\end{tabularx}
\fuente{\archivo{resultados/escenarios/tablas/sensibilidades_plan_recomendado.csv};
referencia central a partir del piloto y los $z$ seleccionados}
\end{table}

El ejemplo del arancel muestra una incertidumbre de supuestos que no aparece
en la banda aleatoria del escenario base: con otro recargo, el pesimista puede
tener contribución negativa. La alternativa sin freno económico mejora el
servicio, pero pierde dinero. Es un conflicto de objetivos que debe comunicarse.

Las pruebas también cubren llegadas iniciales, edades, VM de CSUV viejo,
correlaciones, efecto de producto, dispersión de CSUV, sesgo residual y nivel
de continuidad A014. No se les asignan probabilidades.

\subsection{Por qué se prolonga el análisis más allá de diciembre}
Los pedidos de octubre--diciembre de A013 llegan en enero--marzo de A014.
No seguirlos puede premiar una política que mejora artificialmente el cierre y
desabastece el trimestre siguiente.

La prolongación conserva el nivel desestacionalizado del último trimestre y el
perfil mensual de mercado A009--A010. No repite el adelanto o el lanzamiento en
cada enero. Al no conocerse el calendario ni la economía A014, se presenta como
sensibilidad de continuidad y se compara también un nivel un 10\,\% menor o mayor.

\tabla{continuidad.tex}{Comparación del plan base y recomendado al extender
hasta junio de A014. Resultado medio tras costes modelizados.}
{tab:continuidad}{\archivo{resultados/escenarios/tablas/comparacion_continuidad_politicas.csv}}

La mejora frente al plan base se mantiene bajo esa continuidad: el central
recomendado acumula 93,00 M\euro{} en 18 meses frente a 92,66 M\euro{} del base.
No se añade un precio de venta ficticio por todo el stock que queda ni se
demuestra con ello un óptimo de horizonte infinito.

\subsection{La recomendación operativa, expresada correctamente}
\begin{enumerate}
\item \textbf{Partir del plan central y los colchones seleccionados.} El plan
inicial asigna \num{16144} coches a pedidos decididos en A013 y prevé
\num{18114} llegadas en A013; son flujos distintos. Los anexos muestran mes de
decisión y mes de recepción.
\item \textbf{Revisar mensualmente sólo lo modificable.} Los pedidos realmente
ejecutados después de observar ventas no tienen por qué sumar lo mismo que el
plan inicial. En el central su media es de unos \num{15756} coches.
\item \textbf{Vigilar segmentos, no sólo el total.} Priorizar la disponibilidad
de los segmentos de mayor contribución; vigilar especialmente la incertidumbre
de CSUV, el stock de C y la fragilidad económica de B.
\item \textbf{Separar señales comerciales y restricciones físicas.} Un descenso
de ventas con la red agotada no demuestra menor demanda. Un nuevo coste confirmado
puede justificar frenar pedidos aunque exista demanda.
\item \textbf{Revisar las hipótesis que dominan el resultado.} Disponibilidad de
transfers, aranceles, prima CO$_2$, fechas de producto y tamaño del adelanto no
deben quedar escondidos en una cifra final.
\end{enumerate}

\tabla{revision.tex}{Umbrales orientativos del cociente de sumas de tres meses:
ventas observadas / previsión. Son referencias de seguimiento, no garantías.}
{tab:revision}{\archivo{resultados/escenarios/tablas/umbrales_revision_a013.csv}}

Los percentiles se calculan sobre el mismo cociente utilizado en la regla, sin
unir años de validación distintos. CSUV utiliza un proxy de SSUV/MSUV ampliado,
no un umbral propio validado. El rango muy amplio de S y el rango inferior a uno
de MD también recuerdan los errores históricos y la muestra limitada: las señales
requieren diagnóstico, no una orden automática de aumentar o reducir producción.

\begin{idea}{Respuesta global a las cinco preguntas del caso}
Se observa crecimiento de clientes en el total de A012, distinto de la evolución
de la gama comparable. Se obtiene una previsión A013 evaluada retrospectivamente,
con una limitación específica para CSUV. Los escenarios muestran rango de ventas,
stock y contribución. El plan separa pedidos de llegadas y revisa a tres meses.
La política recomendada cumple el servicio medio central acordado y mejora la
contribución frente al base, con supuestos y riesgos explícitos.
\end{idea}
\FloatBarrier

\appendix
\section{Planes mensuales: pedido y llegada no son intercambiables}\label{anexo:planes}

La primera tabla recoge \textbf{pedidos del plan inicial decididos durante A013}.
Enero corresponde a llegada en abril; diciembre, a marzo de A014. Son cantidades
enteras de una propuesta inicial, no medias ni medianas de decisiones revisadas.
La segunda tabla recoge las \textbf{llegadas iniciales previstas en A013}.
Enero y febrero pertenecen a compromisos anteriores estimados.

\begin{table}[htbp]\centering
\caption{Pedidos iniciales por mes de decisión A013, coches.}\label{tab:pedidos}
{\small\setlength{\tabcolsep}{3pt}\begin{tabular}{lrrrrrrrrrrrr}
\toprule
Segmento & 1 & 2 & 3 & 4 & 5 & 6 & 7 & 8 & 9 & 10 & 11 & 12 \\
\midrule
B & 153 & 132 & 168 & 121 & 66 & 45 & 67 & 78 & 100 & 64 & 103 & 136 \\
C & 0 & 0 & 0 & 0 & 0 & 0 & 0 & 172 & 220 & 163 & 201 & 232 \\
MD & 0 & 15 & 59 & 36 & 17 & 29 & 42 & 33 & 33 & 31 & 40 & 45 \\
SSUV & 659 & 494 & 750 & 530 & 253 & 245 & 325 & 412 & 496 & 219 & 335 & 426 \\
CSUV & 0 & 218 & 347 & 310 & 149 & 112 & 191 & 253 & 272 & 122 & 185 & 257 \\
MSUV & 353 & 622 & 693 & 593 & 293 & 336 & 452 & 516 & 633 & 314 & 309 & 468 \\
S & 80 & 72 & 74 & 30 & 0 & 0 & 0 & 0 & 0 & 56 & 83 & 6 \\
Total & 1.245 & 1.553 & 2.091 & 1.620 & 778 & 767 & 1.077 & 1.464 & 1.754 & 969 & 1.256 & 1.570 \\
\bottomrule
\end{tabular}
}
\fuente{\archivo{resultados/escenarios/tablas/plan_pedidos_inicial_a013_a014.csv};
filtro por año de decisión A013}
\end{table}

\begin{table}[htbp]\centering
\caption{Llegadas del plan inicial por mes de recepción A013, coches.}\label{tab:llegadas}
{\small\setlength{\tabcolsep}{3pt}\begin{tabular}{lrrrrrrrrrrrr}
\toprule
Segmento & 1 & 2 & 3 & 4 & 5 & 6 & 7 & 8 & 9 & 10 & 11 & 12 \\
\midrule
B & 124 & 124 & 81 & 153 & 132 & 168 & 121 & 66 & 45 & 67 & 78 & 100 \\
C & 454 & 454 & 0 & 0 & 0 & 0 & 0 & 0 & 0 & 0 & 172 & 220 \\
MD & 43 & 43 & 0 & 0 & 15 & 59 & 36 & 17 & 29 & 42 & 33 & 33 \\
SSUV & 390 & 390 & 1.135 & 659 & 494 & 750 & 530 & 253 & 245 & 325 & 412 & 496 \\
CSUV & 489 & 489 & 0 & 0 & 218 & 347 & 310 & 149 & 112 & 191 & 253 & 272 \\
MSUV & 576 & 576 & 0 & 353 & 622 & 693 & 593 & 293 & 336 & 452 & 516 & 633 \\
S & 34 & 34 & 329 & 80 & 72 & 74 & 30 & 0 & 0 & 0 & 0 & 0 \\
Total & 2.110 & 2.110 & 1.545 & 1.245 & 1.553 & 2.091 & 1.620 & 778 & 767 & 1.077 & 1.464 & 1.754 \\
\bottomrule
\end{tabular}
}
\fuente{\archivo{resultados/escenarios/tablas/plan_llegadas_inicial_a013.csv}}
\end{table}

Por ejemplo, los \num{1245} coches pedidos en enero de A013 aparecen como
llegadas de abril. Los \num{2110} previstos para enero ya vienen en camino
y no aparecen como un pedido que se pueda decidir en ese mismo enero.
La diferencia anual no es un fallo de balance: cada tabla utiliza un criterio
temporal diferente.
\FloatBarrier

\clearpage
\begin{landscape}
\pagestyle{plain}
\section{Tabla de resultados completa por segmento y escenario}\label{anexo:tabla5}
\noindent\small Medias de 1.000 réplicas. Los totales se calculan por réplica
antes del resumen; cobertura y servicio no son sumas de ratios. El redondeo puede
producir diferencias de una unidad. Los intereses de concesionarios se informan
aparte de la contribución de la marca.
\fuente{\archivo{resultados/escenarios/tablas/tabla5_plan_recomendado_a013.csv}}
{\footnotesize\renewcommand{\arraystretch}{1.02}\setlength{\tabcolsep}{5pt}\begin{longtable}{@{}llrrrrrrrr@{}}
\caption{Volúmenes A013: medias; cantidades en coches salvo cobertura.}\label{tab:anexo-volumenes}\\
\toprule
Seg. & Escenario & \makecell{Demanda\\(coches)} & \makecell{Ventas\\(coches)} & \makecell{Pedidos\\A013} & \makecell{Llegadas\\A013} & \makecell{Stock\\medio} & \makecell{Cobertura\\(meses)} & \makecell{Stock\\final} & \makecell{Stock $>90$\\al cierre} \\
\midrule
\endfirsthead
\toprule
Seg. & Escenario & \makecell{Demanda\\(coches)} & \makecell{Ventas\\(coches)} & \makecell{Pedidos\\A013} & \makecell{Llegadas\\A013} & \makecell{Stock\\medio} & \makecell{Cobertura\\(meses)} & \makecell{Stock\\final} & \makecell{Stock $>90$\\al cierre} \\
\midrule
\endhead
\midrule
\multicolumn{10}{r}{\footnotesize Continúa en la página siguiente.}\\
\endfoot
\bottomrule
\endlastfoot
B & Pesimista & 1.040 & 426 & 0 & 329 & 54,2 & 0,7 & 0 & 0 \\
C & Pesimista & 1.806 & 1.796 & 0 & 908 & 1.072,7 & 7,2 & 219 & 219 \\
MD & Pesimista & 360 & 360 & 286 & 283 & 90,9 & 3,1 & 48 & 5 \\
SSUV & Pesimista & 4.720 & 1.915 & 0 & 1.915 & 169,3 & 0,4 & 0 & 0 \\
CSUV & Pesimista & 2.838 & 2.816 & 1.751 & 2.325 & 979,5 & 4,3 & 520 & 245 \\
MSUV & Pesimista & 4.567 & 4.556 & 4.335 & 4.692 & 936,3 & 2,5 & 672 & 17 \\
S & Pesimista & 322 & 282 & 0 & 397 & 186,4 & 9,3 & 116 & 116 \\
\rowcolor{fondo}
Total & Pesimista & 15.652 & 12.151 & 6.371 & 10.848 & 3.489,4 & 2,7 & 1.574 & 601 \\
\addlinespace[5pt]
B & Central & 1.261 & 1.212 & 1.170 & 1.212 & 180,6 & 1,8 & 97 & 8 \\
C & Central & 2.267 & 2.258 & 959 & 1.298 & 985,2 & 5,3 & 147 & 11 \\
MD & Central & 435 & 433 & 375 & 349 & 83,7 & 2,3 & 41 & 1 \\
SSUV & Central & 5.668 & 5.387 & 4.882 & 5.867 & 700,8 & 1,5 & 480 & 22 \\
CSUV & Central & 3.394 & 3.327 & 2.357 & 2.791 & 919,8 & 3,4 & 475 & 138 \\
MSUV & Central & 5.601 & 5.554 & 5.501 & 5.588 & 862,6 & 1,9 & 570 & 5 \\
S & Central & 390 & 368 & 513 & 778 & 342,6 & 12,2 & 411 & 363 \\
\rowcolor{fondo}
Total & Central & 19.017 & 18.539 & 15.756 & 17.883 & 4.075,3 & 2,6 & 2.221 & 547 \\
\addlinespace[5pt]
B & Optimista & 1.930 & 1.578 & 1.602 & 1.540 & 136,8 & 0,9 & 58 & 1 \\
C & Optimista & 3.633 & 3.451 & 2.283 & 2.403 & 773,2 & 2,6 & 59 & 0 \\
MD & Optimista & 619 & 571 & 546 & 471 & 70,0 & 1,4 & 25 & 0 \\
SSUV & Optimista & 8.560 & 7.077 & 6.822 & 7.388 & 524,4 & 0,8 & 311 & 0 \\
CSUV & Optimista & 5.173 & 4.692 & 3.935 & 4.080 & 804,3 & 1,9 & 399 & 42 \\
MSUV & Optimista & 7.975 & 7.067 & 7.292 & 6.886 & 702,8 & 1,1 & 355 & 0 \\
S & Optimista & 556 & 516 & 788 & 964 & 346,3 & 8,7 & 449 & 361 \\
\rowcolor{fondo}
Total & Optimista & 28.445 & 24.952 & 23.268 & 23.731 & 3.357,8 & 1,4 & 1.656 & 404 \\
\addlinespace[5pt]
\end{longtable}
\clearpage
\begin{longtable}{@{}llrrrrrrrrrr@{}}
\caption{Economía A013: medias; importes en M\euro{} salvo cantidades de pérdidas y transfers.}\label{tab:anexo-economia}\\
\toprule
Seg. & Escenario & \makecell{Pérdidas\\(coches)} & \makecell{Transfers\\(coches)} & \makecell{Coste\\transfer} & \makecell{Coste\\stock} & VM & CVP & CO$_2$ & Arancel & Cambio & Resultado \\
\midrule
\endfirsthead
\toprule
Seg. & Escenario & \makecell{Pérdidas\\(coches)} & \makecell{Transfers\\(coches)} & \makecell{Coste\\transfer} & \makecell{Coste\\stock} & VM & CVP & CO$_2$ & Arancel & Cambio & Resultado \\
\midrule
\endhead
\midrule
\multicolumn{12}{r}{\footnotesize Continúa en la página siguiente.}\\
\endfoot
\bottomrule
\endlastfoot
B & Pesimista & 614 & 0 & 0,00 & 0,04 & 0,73 & 0,05 & 1,09 & 0,16 & 0,00 & -1,24 \\
C & Pesimista & 9 & 0 & 0,00 & 0,52 & 5,39 & 1,21 & 3,51 & 0,34 & 0,00 & -3,16 \\
MD & Pesimista & 0 & 0 & 0,00 & 0,06 & 1,34 & 1,38 & 0,75 & 0,12 & 0,00 & 0,45 \\
SSUV & Pesimista & 2.805 & 0 & 0,00 & 0,12 & 3,79 & 5,47 & 5,42 & 0,96 & 0,00 & -1,03 \\
CSUV & Pesimista & 22 & 0 & 0,00 & 0,68 & 6,90 & 10,48 & 5,38 & 0,90 & 0,00 & 3,52 \\
MSUV & Pesimista & 11 & 0 & 0,00 & 0,68 & 8,89 & 26,56 & 19,15 & 2,01 & 0,00 & 4,72 \\
S & Pesimista & 40 & 0 & 0,00 & 0,10 & 0,59 & 1,34 & 1,27 & 0,14 & 0,00 & -0,16 \\
\rowcolor{fondo}
Total & Pesimista & 3.501 & 0 & 0,00 & 2,20 & 27,63 & 46,49 & 36,56 & 4,64 & 0,00 & 3,09 \\
\addlinespace[5pt]
B & Central & 50 & 0 & 0,00 & 0,13 & 2,08 & 0,15 & 0,00 & 0,00 & 0,00 & 0,02 \\
C & Central & 9 & 0 & 0,00 & 0,55 & 6,60 & 1,69 & 0,00 & 0,00 & 0,00 & 1,15 \\
MD & Central & 2 & 0 & 0,00 & 0,06 & 1,59 & 1,69 & 0,00 & 0,00 & 0,00 & 1,63 \\
SSUV & Central & 281 & 0 & 0,00 & 0,51 & 10,66 & 15,37 & 0,00 & 0,00 & 0,00 & 14,86 \\
CSUV & Central & 67 & 0 & 0,00 & 0,66 & 8,11 & 12,43 & 0,00 & 0,00 & 0,00 & 11,77 \\
MSUV & Central & 47 & 0 & 0,00 & 0,63 & 10,68 & 32,54 & 0,00 & 0,00 & 0,00 & 31,91 \\
S & Central & 23 & 0 & 0,00 & 0,20 & 0,77 & 1,75 & 0,00 & 0,00 & 0,00 & 1,55 \\
\rowcolor{fondo}
Total & Central & 478 & 0 & 0,00 & 2,73 & 40,50 & 65,62 & 0,00 & 0,00 & 0,00 & 62,88 \\
\addlinespace[5pt]
B & Optimista & 352 & 0 & 0,00 & 0,10 & 2,71 & 0,20 & 0,00 & 0,00 & -0,37 & 0,47 \\
C & Optimista & 182 & 0 & 0,00 & 0,53 & 8,73 & 3,94 & 0,00 & 0,00 & -0,59 & 4,00 \\
MD & Optimista & 48 & 0 & 0,00 & 0,05 & 2,07 & 2,25 & 0,00 & 0,00 & -0,11 & 2,31 \\
SSUV & Optimista & 1.483 & 0 & 0,00 & 0,38 & 13,99 & 20,22 & 0,00 & 0,00 & -1,77 & 21,60 \\
CSUV & Optimista & 481 & 0 & 0,00 & 0,59 & 11,29 & 17,68 & 0,00 & 0,00 & -0,92 & 18,01 \\
MSUV & Optimista & 907 & 0 & 0,00 & 0,51 & 13,51 & 41,49 & 0,00 & 0,00 & -1,63 & 42,61 \\
S & Optimista & 40 & 0 & 0,00 & 0,23 & 1,07 & 2,47 & 0,00 & 0,00 & -0,13 & 2,37 \\
\rowcolor{fondo}
Total & Optimista & 3.493 & 0 & 0,00 & 2,38 & 53,36 & 88,23 & 0,00 & 0,00 & -5,52 & 91,37 \\
\addlinespace[5pt]
\end{longtable}
}
\end{landscape}
\pagestyle{fancy}
\clearpage

\section{Cómo reproducir y comprobar el trabajo}\label{anexo:reproducir}

\subsection{La cadena de ejecución}
Desde la raíz del repositorio se ejecutan, con un kernel que tenga las dependencias
instaladas, los cuatro notebooks en orden:
\begin{center}
\archivo{datos.ipynb} $\longrightarrow$ \archivo{demanda.ipynb}
$\longrightarrow$ \archivo{simulacion.ipynb}
$\longrightarrow$ \archivo{escenarios.ipynb}.
\end{center}
Cada uno conserva explicación, código y resultados. Las definiciones compartidas
son celdas \archivo{api}; no se requiere ejecutar manualmente celdas fuera de orden.
Las tablas y figuras se guardan en la carpeta de resultados de cada notebook.
La última fase regenera también el resumen y la presentación.

La ejecución de referencia utiliza semilla 42, 200 réplicas de piloto, 1.000
finales por escenario y 18 meses simulados. Las versiones registradas incluyen
Python 3.13.15, NumPy 2.5.3 y pandas 3.0.6; el registro completo está en
\archivo{resultados/escenarios/tablas/configuracion_ejecucion.json}.

\subsection{Qué verifican las comprobaciones}
\noindent\begin{tabularx}{\linewidth}{@{}p{3.1cm}Y@{}}\toprule
Bloque & Comprobaciones relevantes \\\midrule
Datos & Claves, fechas, lecturas de ausentes, identidad de clientes netos,
balances originales y ajuste de episodios aplicado una sola vez.\\
Previsión & Mercado desde A001; objetivo observado; CSUV fuera de validación;
mutar economía futura no altera la previsión operativa; mutar datos posteriores
no cambia el efecto de producto estimado al origen.\\
Simulador & Ejemplo de 286 coches; FIFO; coste de 15 días; límite de 180;
intereses por tramos; transfers con VM; unidades enteras y balances.\\
Integración & Mismo plan inicial; ausencia de anticipación; pedido de diciembre
a marzo; identidades económicas; totales por réplica y ledger absoluto.\\
Entrega & Coincidencia con tablas; documentos regenerables; legibilidad y
estructura de Word y PowerPoint.\\\bottomrule
\end{tabularx}

Comprobar una identidad contable no demuestra que la hipótesis de demanda sea
correcta. Por eso conviven controles de código, evaluación histórica y
sensibilidades metodológicas: cada uno responde a un tipo diferente de error.

\subsection{Qué representa este informe}
Las tablas y figuras de \archivo{docs/report/} son una instantánea de la ejecución
citada. Compilar LaTeX no vuelve a estimar los modelos. Las figuras vectoriales
didácticas se dibujan en el propio documento a partir de reglas o de copias de
resultados; no crean una segunda simulación. Si cambia el análisis, deben
actualizarse también las cifras y las interpretaciones del informe.

\section{Preguntas de repaso y errores de interpretación frecuentes}\label{anexo:repaso}

\begin{enumerate}[itemsep=7pt]
\item \textbf{¿Por qué no son iguales matriculaciones y demanda de clientes?}
Porque las primeras incluyen automatriculaciones. Además, ni siquiera las compras
netas identifican toda la demanda potencial cuando falta stock.
\item \textbf{¿Puede crecer el total y caer la gama comparable?}
Sí. CSUV incorpora clientes nuevos al perímetro de la marca aunque los segmentos
anteriores permanezcan estables o retrocedan.
\item \textbf{¿Por qué no se descuenta dos veces el adelanto?}
El histórico ajustado sirve para interpretar o normalizar el ancla; el registro
futuro se aplica una sola vez a A013. Replicar el descuento anual en cada mes
o añadir otro efecto de caída sobre la misma cantidad duplicaría su impacto.
\item \textbf{¿Un retraso económico de tres meses permite conocer todo A013?}
No. Sólo parte del horizonte queda cubierta por información ya observada; para
el resto se necesita congelación o escenarios.
\item \textbf{¿Qué significa MAE relativo del 14,8\,\%?}
Es el error absoluto mensual del total comparable en la evaluación retrospectiva
A012, dividido por sus ventas observadas. No valida CSUV ni garantiza el error A013.
\item \textbf{¿La cifra prevista es un pedido a fábrica?}
No. Antes se descuentan inventario y compromisos, se añade el colchón y se respeta
el plazo de llegada.
\item \textbf{¿95\,\% de servicio significa un 95\,\% de probabilidad de no tener roturas?}
No. El objetivo es la media de la proporción anual de unidades atendidas en el
central. Un segmento puede cumplirla y tener réplicas por debajo del umbral.
\item \textbf{¿Las medianas de segmentos suman la mediana del total?}
No necesariamente. Hay que sumar segmentos dentro de cada réplica y calcular
después el percentil del total.
\item \textbf{¿Toda venta perdida es una pérdida económica positiva?}
No. La venta adicional también habría soportado sus costes. Si la contribución
marginal es negativa, perder esa operación no equivale a perder un margen positivo.
\item \textbf{¿Más réplicas resuelven la incertidumbre de un arancel supuesto?}
No. Reducen ruido de Monte Carlo bajo ese arancel. Para estudiar otro tamaño
hay que cambiar el supuesto y comparar una sensibilidad.
\item \textbf{¿La mejor política evaluada es el óptimo global?}
No. Es la mejor dentro de la malla y del objetivo acordados, bajo las restricciones
y aproximaciones representadas.
\item \textbf{¿Qué dato vigilarías primero en A013?}
Ventas respecto a previsión, disponibilidad efectiva y compromisos; junto con
confirmación de costes, lanzamientos e interrupciones. Ninguna señal debe
interpretarse sin las restricciones que pueden estar ocultando demanda.
\end{enumerate}


\clearpage
\begin{thebibliography}{9}
\addcontentsline{toc}{section}{Referencias y trazabilidad}
\bibitem{caso} J. Velasco y J. Goikoetxea.
\emph{Kaido Motors España: demanda, stock y beneficio}.
Universidad Alfonso X el Sabio. Enunciado docente, 16 páginas.
\archivo{docs/kaido-motors-case-study.pdf}.
\bibitem{guia} Proyecto Kaido Motors.
\emph{Guía del proyecto y plan aprobado de revisión del modelo}.
\archivo{docs/guia-proyecto.md} y \archivo{docs/plan-revision-modelo.md}.
\bibitem{datos} Proyecto Kaido Motors.
\emph{Datos y demanda observada}. Notebook \archivo{datos.ipynb};
salidas en \archivo{resultados/datos/}.
\bibitem{demanda} Proyecto Kaido Motors.
\emph{Previsión de clientes y escenarios A013}. Notebook \archivo{demanda.ipynb};
salidas en \archivo{resultados/demanda/}.
\bibitem{simulacion} Proyecto Kaido Motors.
\emph{Simulación logística A013}. Notebook \archivo{simulacion.ipynb};
salidas en \archivo{resultados/simulacion/}.
\bibitem{escenarios} Proyecto Kaido Motors.
\emph{Escenarios y recomendación A013}. Notebook \archivo{escenarios.ipynb};
salidas en \archivo{resultados/escenarios/}.
\bibitem{configuracion} Proyecto Kaido Motors.
\emph{Configuración de la ejecución revisada}.
\archivo{resultados/escenarios/tablas/configuracion_ejecucion.json}.
\end{thebibliography}
\end{document}
